% 图 3 系统运行 — 基于 CICC0903540 技术文档 3
\documentclass[tikz,border=8pt]{standalone}
% 顶会/顶刊风格：低饱和度马卡龙 + 高级感
% 适用于 NeurIPS, CVPR, ICLR, IEEE TPAMI
\usepackage{amssymb}
\usepackage{fontspec}
\usepackage{xeCJK}
\setCJKmainfont{SimHei}
\usepackage{tikz}
\usepackage{pgfplots}
\usetikzlibrary{
  arrows.meta, positioning, shapes.geometric, calc,
  backgrounds, fit, decorations.pathreplacing, patterns
}

% 低饱和度马卡龙配色
\definecolor{mBlue}{HTML}{AEC6CF}
\definecolor{mPink}{HTML}{FFD1DC}
\definecolor{mGreen}{HTML}{B1D8B7}
\definecolor{mPurple}{HTML}{B39EB5}
\definecolor{mYellow}{HTML}{FDFD96}
\definecolor{mGray}{HTML}{E5E4E2}
\definecolor{mDark}{HTML}{4A4A4A}

% 全局 TikZ 样式
\tikzset{
  >=stealth,
  every node/.style={align=center, font=\sffamily},
  block/.style={
    rounded corners=3pt, draw=mDark, align=center,
    minimum height=0.8cm, inner sep=5pt},
  arr/.style={-{Stealth[length=4pt]}, semithick, draw=mDark},
  % 信息密度增强：张量维度标注
  ts/.style={font=\tiny\sffamily, color=mDark!80},
  % 数学操作符节点（⊕加 ⊗乘 ‖拼接）
  op/.style={circle, draw=mDark, minimum size=0.4cm, inner sep=0,
    font=\scriptsize\sffamily, fill=mGray!60},
  % 图例
  legendbox/.style={rectangle, draw=mDark!50, rounded corners=2pt,
    fill=mGray!15, inner sep=4pt, font=\tiny\sffamily},
}

\begin{document}
\begin{tikzpicture}[
  node distance=0.5cm and 1.2cm,
  block/.style={
    rectangle, rounded corners=3pt, draw=mDark, align=center,
    font=\small\sffamily, semithick,
    minimum width=2.5cm, minimum height=0.75cm, inner sep=4pt},
  sub/.style={
    rectangle, rounded corners=2pt, draw=mDark!60, align=center,
    font=\scriptsize\sffamily, color=mDark,
    minimum width=2.0cm, minimum height=0.5cm, inner sep=3pt},
  arr/.style={-{Stealth[length=4pt]}, draw=mDark, semithick},
  sect/.style={font=\footnotesize\sffamily\bfseries, color=mDark},
]
  \node[sect] (sectHost) {上位机};
  \node[block, fill=mBlue!40, below=0.35cm of sectHost] (load)
    {模型加载\\{\footnotesize test\_sub.py}};
  \node[block, fill=mBlue!40, below=0.4cm of load] (enc)
    {Encoder 推理\\{\footnotesize 语义特征提取}};
  \node[block, fill=mBlue!40, below=0.4cm of enc] (quant)
    {量化打包\\{\footnotesize scale, zero\_point}};
  \node[sub, fill=mGray!40, below=0.2cm of quant] (sshInit)
    {SSHClient, AutoAddPolicy, compress};

  \node[block, fill=mGreen!50, right=2.6cm of quant] (trans)
    {OpenSSH 传输\\{\footnotesize SFTP 连接池}};
  \node[sub, fill=mGray!40, below=0.2cm of trans] (sftpOpt)
    {aes128-ctr, 16--64MB 分块, mmap};

  \node[block, fill=mPurple!40, right=2.6cm of trans] (recv)
    {接收 latent\\{\footnotesize .pt 文件}};
  \node[block, fill=mPurple!40, below=0.4cm of recv] (dequant)
    {去量化 + AWGN};
  \node[block, fill=mPurple!40, below=0.4cm of dequant] (gen)
    {Generator 推理\\{\footnotesize TVM / MNN}};
  \node[block, fill=mPurple!40, below=0.4cm of gen] (save)
    {PNG 落盘};
  \node[sect, above=0.35cm of recv] (sectPi) {飞腾派};

  \draw[arr] (load) -- (enc);
  \draw[arr] (enc) -- (quant)
    node[pos=0.5, left=2pt, ts] {$1{\times}32{\times}32{\times}32$};
  \draw[arr] (quant.east) -- (trans.west)
    node[pos=0.5, above=4pt, ts] {语义张量};
  \draw[arr] (trans.east) -- (recv.west)
    node[pos=0.5, above=4pt, ts] {$\hat{\mathbf{y}}$};
  \draw[arr] (recv) -- (dequant);
  \draw[arr] (dequant) -- (gen);
  \draw[arr] (gen) -- (save)
    node[pos=0.5, left=2pt, ts] {PNG};

  \begin{scope}[on background layer]
    \node[fit=(load)(enc)(quant)(sshInit), fill=mBlue!12, draw=mDark!25,
      rounded corners=4pt, inner sep=5pt, semithick] {};
    \node[fit=(recv)(dequant)(gen)(save), fill=mPurple!12, draw=mDark!25,
      rounded corners=4pt, inner sep=5pt, semithick] {};
  \end{scope}

  % 图例（预留边距，避免与飞腾派列重叠）
  \node[legendbox, anchor=north east] at ([xshift=-5pt, yshift=-5pt]current bounding box.south east)
    [align=left, text width=3cm] {
    \textbf{图例}\\[1pt]
    \colorbox{mBlue!40}{\rule{0.15cm}{0.1cm}\hspace{0.06cm}} 上位机 \quad
    \colorbox{mPurple!40}{\rule{0.15cm}{0.1cm}\hspace{0.06cm}} 飞腾派 \quad
    \colorbox{mGreen!50}{\rule{0.15cm}{0.1cm}\hspace{0.06cm}} 传输
  };
\end{tikzpicture}
\end{document}
